% Cote 119: offprints, reports, courses, CV and letters (typed and printed part): archive cote 119, site page 300 (PDF page 301).
% typescript. Draft compared: none.
% AI-transcribed by Claude Opus 5.5 (claude-opus-5-5, adaptive thinking, effort high), Message Batches API, batch r119_full (msgbatch_01DwZNm5Y7YsKBWyAN4sFs8w), 2026-11-13; not yet checked by eye.
\scanpage{119/300}
\mspage{- 13 -}

SÉMINAIRE

B.KAHN - \uline{Formes de Pfister et invariants cohomologiques} (2 exposés).

E.BAYER-FLUCKIGER - \uline{Le théorème de simplification dans le cas hermitien}.

SÉMINAIRE COMMUN AVEC LA CHAIRE DE THÉORIE DES GROUPES

J-P.SERRE - \uline{Travaux de Chernousov sur les groupes de type} $E_8$ .

J.TITS - \uline{Travaux de Chernousov sur les groupes de type} $E_8$ (2 exposés).

J-P.SERRE - \uline{Remarques sur la cohomologie galoisienne des groupes semi-simples}.

PUBLICATIONS

J-P.SERRE, \uline{Construction de revêtements étales de la droite affine en}
\uline{caractéristique} $p$ (\uline{C.R.Acad.Sci.Paris}, 311, 1990, série I, 341-346).

-- \uline{Spécialisation des éléments de} $\mathrm{Br}_2(\mathbf{Q}(T_1,\dots,T_n))$ (\uline{C.R.Acad.Sci.Paris},
311, 1990, série I, 397-402).

-- \uline{Relèvements dans} $\tilde{A}_n$ (\uline{C.R.Acad.Sci.Paris}, 311, 1990, série I, 477-482).

-- \uline{Revêtements à ramification impaire et thêta-caractéristiques} (\uline{C.R.Acad}.
\uline{Sci.Paris}, 311, 1990, série I, 547-552).

-- \uline{Les petits cousins} (\uline{Miscellanea Mathematica}, Springer-Verlag, 1991,
277-291).

MISSIONS

Cours

- \uline{Topics in Galois cohomology}, Harvard, septembre-décembre 1990;

- \uline{Sieves}, Singapour, mai 1991.

Exposés

- \uline{How often does a conic have a rational point} ? , State College, septembre
1990; Yale, novembre 1990;

- \uline{Coverings of algebraic curves}, Harvard, octobre 1990;

- \uline{Motives}, Harvard, octobre 1990;

- \uline{Riemann Hypothesis} : \uline{Why} ? , Chicago, octobre 1990;

- \uline{Galois groups of division points of abelian varieties}, Chicago, octobre
1990;

- \uline{Bounds for number of points of hypersurfaces over finite fields}, Chicago,
octobre 1990;

\mspage{\hand{300}}
